% Image credits for the photographs and AI illustrations of Book 5 (University, Year 3).
% Machine-readable license records live in images/book5/CREDITS.md; keep the
% two files in sync. The Book 5 writer owns this file: add one \item per
% photograph (inside the itemize, by chapter) as photographs land.
% Portuguese (Brazilian) edition: translated from frontmatter/image-credits-book5.tex;
% every credit, name, dataset and licence kept exactly.
\chapter*{Créditos das imagens}

Os desenhos, os mapas e os gráficos deste livro são originais; os mapas e os
gráficos desenhados a partir de dados publicados citam sua fonte na legenda ou
nesta lista. As fotografias são usadas sob suas respectivas licenças livres, com
agradecimentos a seus autores; os links completos para cada fotografia e as URLs
das licenças estão listados em \texttt{images/book5/CREDITS.md}, no repositório
do livro.

\bigskip
Ao lado das fotografias e dos desenhos originais, algumas figuras são ilustrações
geradas por inteligência artificial, produzidas com a geração de imagens da OpenAI
a partir de prompts escritos pelos editores do livro e revisadas quanto à precisão
científica antes de sua inclusão. Elas mostram apenas cenas genéricas, nunca um
lugar real, uma pessoa real ou um espécime real. O prompt completo de cada imagem
gerada está registrado em \texttt{images/book5/ai/PROMPTS.md}, no repositório.

\bigskip
\noindent Fotografias e dados, por capítulo:
\begin{itemize}\small
  \item Cap.~1 (sismologia e imageamento da Terra): PREM, A.~M.~Dziewonski e D.~L.~Anderson (1981), conforme tabulado no ObsPy; ak135, B.~L.~N.~Kennett, E.~R.~Engdahl e R.~Buland (1995).
  \item Cap.~3 (gravidade e a figura da Terra): os satélites gêmeos GRACE Follow-On, concepção artística, NASA/JPL-Caltech, domínio público (Wikimedia Commons). Balanço de massa dos mantos de gelo da equipe IMBIE (2021), UK Polar Data Centre, Open Government Licence v3.0.
  \item Cap.~4 (geomagnetismo e paleomagnetismo): Campo Geomagnético Internacional de Referência, décima quarta geração, IAGA Division~V Working Group V-MOD, coeficientes distribuídos por NOAA NCEI.
  \item Cap.~5 (modelos térmicos da litosfera): idade do assoalho oceânico, M.~Seton e outros (2020), EarthByte, CC BY 4.0; relevo global ETOPO1, NOAA NCEI, domínio público.
  \item Cap.~6 (convecção no manto e forças motrizes das placas): as praias soerguidas de Ballyhillin, Malin Head, A.~F.~Borchert, CC BY-SA 4.0 (Wikimedia Commons); PREM, A.~M.~Dziewonski e D.~L.~Anderson (1981).
  \item Cap.~7 (mecanismos de deformação e zonas de cisalhamento): milonito com porfiroclastos oculares, Mikenorton, CC BY-SA 4.0 (Wikimedia Commons).
  \item Cap.~9 (tectônica extensional e transcorrente): topografia de bacias e cadeias ao sul do Death Valley, Ken Lund, CC BY-SA 2.0 (Wikimedia Commons).
  \item Cap.~14 (geoquímica de isótopos estáveis): empilhamento bentônico de $\delta^{18}$O LR04, L.~E.~Lisiecki e M.~E.~Raymo (2005).
  \item Cap.~15 (interação água--rocha): terraços de travertino de Palette Spring, Mammoth Hot Springs, Jim Peaco, National Park Service, domínio público (Wikimedia Commons). Constantes de equilíbrio da base de dados do PHREEQC, phreeqc.dat, D.~L.~Parkhurst e C.~A.~J.~Appelo, U.S. Geological Survey, domínio público.
  \item Cap.~16 (estratigrafia de sequências): os Book Cliffs acima de Helper, Utah, David Mackie, domínio público (Wikimedia Commons); empilhamento bentônico de $\delta^{18}$O LR04, L.~E.~Lisiecki e M.~E.~Raymo (2005).
  \item Cap.~17 (modelos de fácies): estratificação cruzada \emph{hummocky} na Formação Arisaig, Michael C.~Rygel, CC BY-SA 3.0 (Wikimedia Commons), recortada.
  \item Cap.~18 (geomorfologia quantitativa): o desfiladeiro do Niágara e as Horseshoe Falls, Ken Lund, CC BY-SA 2.0 (Wikimedia Commons), recortada.
  \item Cap.~19 (dinâmica de fluidos geofísicos): ventos em média zonal da NCEP--NCAR Reanalysis~1, dados fornecidos pelo NOAA PSL, Boulder, Colorado, EUA; World Ocean Atlas 2023, NOAA NCEI.
  \item Cap.~20 (dinâmica atmosférica): um sistema de baixa pressão ao largo da Islândia, NASA/GSFC, MODIS Rapid Response Team, Jacques Descloitres, domínio público (Wikimedia Commons); ventos em média zonal da NCEP--NCAR Reanalysis~1, dados fornecidos pelo NOAA PSL, Boulder, Colorado, EUA.
  \item Cap.~21 (dinâmica dos oceanos): um flutuador Argo lançado de um navio de pesquisa, NOAA Global Ocean Monitoring and Observing Program, domínio público (Wikimedia Commons).
  \item Cap.~22 (modelos climáticos): supercomputadores do NASA Center for Climate Simulation, Goddard Space Flight Center, 2010, domínio público (Wikimedia Commons); temperatura do ar à superfície e radiação no topo da atmosfera da NCEP--NCAR Reanalysis~1, dados fornecidos pelo NOAA PSL, Boulder, Colorado, EUA.
  \item Cap.~24 (glaciologia): uma geleira de descarga no sudoeste da Groenlândia, NASA / Jim Yungel, Operation IceBridge, 2014, domínio público (Wikimedia Commons); altitudes da superfície da Groenlândia do ETOPO1 Ice Surface, NOAA NCEI; extensões do gelo marinho ártico do NSIDC Sea Ice Index, versão~4 (Fetterer e outros, 2025, doi:10.7265/a98x-0f50).
  \item Cap.~25 (sensoriamento remoto e SIG): o interferograma de 1995--2001 do centro vulcânico de Three Sisters, Oregon, U.S. Geological Survey (C.~Wicks e outros, 2002), domínio público (Wikimedia Commons); espectros de laboratório da ECOSTRESS Spectral Library, versão~1.0 (S.~K.~Meerdink, S.~J.~Hook, D.~A.~Roberts e E.~A.~Abbott, 2019; A.~M.~Baldridge e outros, 2009), reproduzidos por cortesia do Jet Propulsion Laboratory, California Institute of Technology, Pasadena, California.
  \item Cap.~26 (interiores e superfícies planetárias): a Cratera do Meteoro de Barringer, Arizona, USGS/\allowbreak D.~Roddy, domínio público (Wikimedia Commons); o Olympus Mons, mosaico do orbitador Viking, NASA/\allowbreak JPL/\allowbreak USGS, domínio público (Wikimedia Commons). Funções de produção e de cronologia de crateras conforme o software Craterstats de G.~G.~Michael, segundo G.~Neukum, B.~A.~Ivanov e W.~K.~Hartmann (2001).
  \item Cap.~27 (atmosferas planetárias e habitabilidade): uma rede de vales secos em Marte, Mars Express HRSC, ESA/DLR/FU Berlin, CC BY-SA 3.0 IGO (Wikimedia Commons); Ligeia Mare, Titã, mosaico de radar da Cassini, NASA/JPL/USGS, domínio público (Wikimedia Commons). Ajustes da zona habitável de R.~K.~Kopparapu e outros (2013); pressão de vapor do metano do NIST Chemistry WebBook.
  \item Cap.~28 (exoplanetas): HIP~65426~b imageado pelo James Webb Space Telescope, NASA/ESA/CSA, A~Carter (UCSC), a equipe ERS~1386 e A.~Pagan (STScI), CC BY~4.0 (Wikimedia Commons). Planetas confirmados do NASA Exoplanet Archive, tabela Planetary Systems Composite Parameters, acesso em 23 de outubro de 2026: esta pesquisa utilizou o NASA Exoplanet Archive, operado pelo California Institute of Technology, sob contrato com a National Aeronautics and Space Administration no âmbito do Exoplanet Exploration Program. Curvas-modelo massa--raio de L.~Zeng e outros (2016, 2019).
  \item Cap.~29 (a Via Láctea e as galáxias): o panorama da Via Láctea do projeto GigaGalaxy Zoom, ESO/S.~Brunier, CC BY~4.0 (Wikimedia Commons); a galáxia espiral M101, Agência Espacial Europeia e NASA (Hubble, heic0602; dados do Hubble K.~D.~Kuntz, F.~Bresolin, J.~Trauger, J.~Mould e Y.-H.~Chu; processamento de imagem Davide De~Martin), CC BY~4.0 (Wikimedia Commons). Velocidades circulares da Galáxia de A.-C.~Eilers, D.~W.~Hogg, H.-W.~Rix e M.~K.~Ness (2019); modelo de massa de E.~Pouliasis, P.~Di~Matteo e M.~Haywood (2017); cefeidas clássicas OGLE da Grande Nuvem de Magalhães, I.~Soszy\'nski e outros (2015), Acta Astronomica 65, 297, distribuídas pelo CDS, Strasbourg (VizieR J/AcA/65/297).
  \item Cap.~30 (de onde vêm os elementos): a Nebulosa do Caranguejo, NASA, ESA, J.~Hester e A.~Loll (Arizona State University), domínio público (Wikimedia Commons); a quilonova de GW170817 in NGC~4993, NASA e ESA, agradecimento a A.~J.~Levan, N.~R.~Tanvir, A.~Fruchter e O.~Fox (Hubble, heic1717a), CC BY~4.0 (Wikimedia Commons). Abundâncias solares de M.~Asplund, A.~M.~Amarsi e N.~Grevesse (2021); seções de choque de captura de nêutrons do KADoNiS v0.3 (I.~Dillmann, R.~Plag, F.~K\"appeler e T.~Rauscher); meias-vidas do IAEA Live Chart of Nuclides (ENSDF); metalicidades do levantamento de Geneva--Copenhagen reanalisadas por L.~Casagrande e outros (2011), distribuídas pelo CDS, Strasbourg (VizieR J/A+A/530/A138).
\end{itemize}
